\begin{trapbox}

	\begin{description}[style=nextline, leftmargin=2.8em, labelsep=0.5em, itemsep=0.35em]
		\item[\textbf{\textcolor{CTrap}{易错点一（忽略有界条件）}}]
			认为不管可行域有无界，最值一定在顶点取到。
		\item[\textbf{\textcolor{CTrap}{正确理解（检查有界性）：}}]
			结论要求“可行域有界”或“最优解存在”；若可行域无界，可能没有有限最值，不能直接套用顶点法。
		\item[\textbf{\textcolor{CTrap}{错因分析（前提遗漏）：}}]
			只死记“顶点取最值”，忽略该结论成立的前提条件。
	\end{description}

	\medskip
	\begin{description}[style=nextline, leftmargin=2.8em, labelsep=0.5em, itemsep=0.35em]
		\item[\textbf{\textcolor{CTrap}{易错点二（截距方向混淆）}}]
			求最值时，直接认为目标函数纵截距越大 $z$ 越大，未考虑 $b$ 的正负。
		\item[\textbf{\textcolor{CTrap}{正确理解（符号决定单调）：}}]
			当 $b>0$ 时，$z$ 与截距同向；$b<0$ 时，$z$ 与截距反向。解题时先判断 $b$ 的符号。
		\item[\textbf{\textcolor{CTrap}{错因分析（直接代公式）：}}]
			记混了符号关系，没有从 $z=ax+by$ 变形为 $y=-\frac{a}{b}x+\frac{z}{b}$ 推导。
	\end{description}

	\medskip
	\begin{description}[style=nextline, leftmargin=2.8em, labelsep=0.5em, itemsep=0.35em]
		\item[\textbf{\textcolor{CTrap}{易错点三（顶点遗漏或算错）}}]
			求可行域顶点时，漏掉某些边角交点，或者解方程组出错，导致最值计算错误。
		\item[\textbf{\textcolor{CTrap}{正确理解（完整枚举）：}}]
			必须将约束条件的所有边界直线两两配对求解，并验证是否满足全部约束，才能得到完整顶点集。
		\item[\textbf{\textcolor{CTrap}{错因分析（计算草率）：}}]
			解方程不细心，或者忽视某些边界（如坐标轴），导致顶点不完整。
	\end{description}

\end{trapbox}
